\section{Motivation}

Städte weltweit veröffentlichen zunehmend große Mengen an offenen Daten \cite{open_gov}.
Wien ist dabei besonders aktiv: Das städtische Open-Data-Portal \textit{data.wien.gv.at}
stellt über 300 Datensätze frei zur Verfügung \cite{datawien}.
Trotz dieser Fülle zeigt unsere Konkurrenzanalyse aus Meilenstein~1 ein ernüchterndes Bild:
Von den insgesamt $n = 345$ auf \textit{data.gv.at} gelisteten Anwendungen für den
Publisher \enquote{Stadt Wien} visualisieren nur rund 100 tatsächlich offene Daten.
Filtert man auf mobile Apps, verbleiben lediglich \textbf{$n = 8$ Anwendungen} –
und selbst diese sind überwiegend hoch spezialisiert oder veraltet.
Im Bereich allgemeiner Open-Data-Visualisierung für Wien besteht somit fast keine
direkte Konkurrenz.

Die Hürde liegt nicht in der Verfügbarkeit der Daten, sondern in ihrer \emph{Zugänglichkeit}.
Simonofski et al.\ \cite{simonofski2022} zeigen, dass gut gestaltete Dashboards das
Verständnis offener Daten bei Bürger:innen deutlich erhöhen – aber nur dann, wenn
heterogene Nutzergruppen von Beginn an berücksichtigt werden.
In der Praxis werden Open-Data-Datensätze jedoch häufig ohne Rücksicht auf
Nutzerinteressen veröffentlicht \cite{open_gov}, in schlecht verarbeitbaren
Formaten und ohne visuelle Aufbereitung.

Dieses Projekt begegnet diesem Problem mit \textbf{MapApp}: einer nativen
Android-Anwendung, die Wiens offene Datensätze auf einer interaktiven Karte
räumlich erkundbar macht.
Statt ein Datenportal zu besuchen und Rohdaten zu interpretieren, sehen
Nutzer:innen Datensatzinhalte als farbkodierte, interaktive Kartenmarkierungen –
ergänzt durch Offline-Downloads, Suchfunktion und Barrierefreiheitsoptionen.

Aus der Nutzeranalyse in Meilenstein~1 ergaben sich drei primäre Nutzergruppen
sowie eine sekundäre:
\begin{itemize}
  \item \textbf{Technikaffine Studierende} (Benjamin Eberhart) – suchen gezielt
        nach spezifischen Datensätzen und wünschen sich detaillierte Filteroptionen.
  \item \textbf{Content Creator und Blogger} (Leah Reiniger) – benötigen visuell
        ansprechende Inhalte, die einfach geteilt werden können.
  \item \textbf{Berufliche Nutzer:innen} (Ralph Ebersbach) – verlassen sich auf
        Datengenauigkeit und brauchen schnellen, wiederholbaren Zugang zu
        bestimmten Datensätzen.
  \item \textbf{Gelegenheitsnutzer:innen} (Luca Kunze) – möchten ohne
        Vorkenntnisse unkompliziert zu interessanten Informationen gelangen.
\end{itemize}

Die zentrale Gestaltungsherausforderung besteht darin, \emph{Informationsdichte}
und \emph{Benutzerfreundlichkeit} in Einklang zu bringen –
eine Spannung, die den gesamten Designprozess über vier Meilensteine geprägt hat.
